\documentclass[11pt]{article}

\usepackage[T1]{fontenc}
\usepackage{lmodern}
\usepackage[margin=1in]{geometry}
\usepackage{amsmath,amssymb,amsthm}
\usepackage{microtype}
\usepackage{xcolor}
\usepackage[colorlinks=true,linkcolor=blue!55!black,urlcolor=blue!55!black]{hyperref}

\newtheorem{proposition}{Proposition}
\newcommand{\dd}{\mathrm{d}}
\setlength{\parindent}{0pt}
\setlength{\parskip}{0.45em}

\title{Analytical Results in the Quantitative Model}
\author{}
\date{}

\begin{document}
\maketitle
\vspace{-2.5em}

Several margins of the quantitative model remain analytical even though its
dynamic program is solved numerically.  They make the housing--fertility
mechanism and demographic closure transparent while retaining the model's
discrete housing products, sequential fertility choices, and full age
distribution.

\section{Stationary demography}

Let $\mathbf p$ collect persons by sex and single year of age.  The matrix $A$
ages and survives the existing population, $\mathbf f$ gives surviving
newborns per birth, $\mathbf d$ contains headship rates, and $\mathbf m$ is
fixed net migration.  If $b$ denotes annual births per household head, the
person law is
\begin{equation}
 \mathbf p'=A\mathbf p+\mathbf f b\,\mathbf d'\mathbf p+\mathbf m.
 \label{eq:person_law}
\end{equation}
The law is linear except for the feedback from the number of household heads
to births.

Define the lifetime population response to one annual birth and the response
to fixed migration as
\begin{equation}
 \mathbf v=(I-A)^{-1}\mathbf f,
 \qquad
 \mathbf w=(I-A)^{-1}\mathbf m.
 \label{eq:linear_responses}
\end{equation}
Two scalars summarize their implications for household formation:
\begin{equation}
 L=\mathbf d'\mathbf v,
 \qquad
 M_H=\mathbf d'\mathbf w,
 \qquad
 R=bL.
 \label{eq:renewal_objects}
\end{equation}
Here $L$ is the lifetime number of household heads generated by one annual
birth, $M_H$ is the stock of heads supported by fixed migration in the absence
of endogenous births, and $R$ is the demographic renewal ratio.

\begin{proposition}[Stationary population]
\label{prop:stationary_population}
Suppose $I-A$ is invertible.  If migration is fixed in levels, $M_H>0$, and
$R<1$, the stationary number of births $B^*$, household heads $S^*$, and
persons $\mathbf p^*$ are
\begin{equation}
 B^*=\frac{bM_H}{1-R},
 \qquad
 S^*=\frac{M_H}{1-R},
 \qquad
 \mathbf p^*=\mathbf w+\frac{bM_H}{1-R}\mathbf v.
 \label{eq:stationary_population}
\end{equation}
The implied vector $\mathbf p^*$ must be nonnegative.  With no migration, a
nonzero stationary population requires $R=1$; the housing market then
determines its scale.
\end{proposition}

\begin{proof}
At a stationary point, $B=bS$ and \eqref{eq:person_law} gives
$\mathbf p=\mathbf w+B\mathbf v$.  Hence
$S=\mathbf d'\mathbf p=M_H+BL$ and
$B=b(M_H+BL)$.  Solving the last equation yields
\eqref{eq:stationary_population}.  When $\mathbf m=0$, a positive fixed point
instead requires $B=bBL$, or $R=1$.
\end{proof}

This proposition is the full age-structured counterpart of a replacement-rate
condition in a two-period model.  The ratio $R$ measures lifetime replacement;
it is distinct from an eigenvalue measuring the speed at which a given age
distribution converges.

\section{Household margins}

The Bellman problem remains numerical, but three conditional margins have
simple forms.

\subsection{Rental housing}

Conditional on next-period assets, let $X$ be resources left after saving and
the consumption and housing floors have been paid.  Write $\alpha$ for the
consumption share in the within-period Stone--Geary aggregator.  When the
rental cap is slack, discretionary consumption and housing satisfy
\begin{equation}
 c-\bar c=\alpha X,
 \qquad
 h-\bar h=\frac{1-\alpha}{r}X.
 \label{eq:renter_solution}
\end{equation}
If the second expression would put total housing above $h_R^{\max}$, then
$h=h_R^{\max}$ and consumption absorbs the remaining resources.  Equivalently,
on a fixed branch the rental condition is
\begin{equation}
 \frac{u_h}{u_c}=r+\zeta^R,
 \qquad
 \zeta^R\geq0,
 \label{eq:rental_wedge}
\end{equation}
where $\zeta^R$ is positive only when the upper rental limit binds.  The
child-space floor raises $\bar h$ and therefore reduces discretionary space; it
is not a second upper-cap wedge.

\subsection{Discrete housing products}

Owner housing is a discrete product in the quantitative model.  Let $V_j$ be
the optimized value of alternative $j$, including its down payment,
transaction cost, and continuation value.  With logit scale $\kappa_T>0$, the
choice probability and its local change are
\begin{align}
 \pi_j&=\frac{\exp(V_j/\kappa_T)}{\sum_k\exp(V_k/\kappa_T)},
 \label{eq:tenure_logit}\\
 \dd\pi_j&=\frac{\pi_j}{\kappa_T}
 \left(\dd V_j-\sum_k\pi_k\dd V_k\right).
 \label{eq:tenure_logit_derivative}
\end{align}
These expressions apply while the feasible set of products is unchanged.
Down-payment constraints and other feasibility changes make the derivative
branch-specific.  Value differences across adjacent products, rather than a
continuous owner-housing first-order condition, are therefore the natural
measure of access in this model.

\subsection{Sequential fertility}

At age $a$, let $V_a^W$ be the value of waiting, $V_a^B$ the continuation
value following conception, $f_a$ the conception probability conditional on
an attempt, and $C_a$ a fixed birth cost.  The value of attempting and its
difference from waiting are
\begin{equation}
 V_a^A=(1-f_a)V_a^W+f_a(V_a^B-C_a),
 \qquad
 \Delta V_a=f_a(V_a^B-C_a-V_a^W).
 \label{eq:attempt_value}
\end{equation}
With fertility-choice scale $\kappa_a$, the attempt probability and realized
birth hazard are
\begin{equation}
 \pi_a^F=\Lambda\!\left(\frac{\Delta V_a}{\kappa_a}\right),
 \qquad
 b_a=f_a\pi_a^F.
 \label{eq:birth_hazard}
\end{equation}
For a smooth shifter $g$ that leaves $f_a$, $C_a$, and $\kappa_a$ fixed,
\begin{equation}
 \frac{\partial b_a}{\partial g}
 =\frac{f_a^2\pi_a^F(1-\pi_a^F)}{\kappa_a}
 \frac{\partial(V_a^B-V_a^W)}{\partial g}.
 \label{eq:birth_hazard_derivative}
\end{equation}
Housing affects births only insofar as it changes the value of the next child
state relative to waiting.  Equation~\eqref{eq:birth_hazard_derivative} also
shows two sources of attenuation: conception risk and a nearly deterministic
attempt margin.

\section{Long-run equilibrium}

Let $z=(\log P,T)$ collect the terminal house price and equal transfer.  Given
a policy or preference shifter $g$, solving the household and distribution
blocks at $(z,g)$ produces a renewal ratio $R(z,g)$, occupied housing per head
$\bar h(z,g)$, and housing supply $K(P,g)$.

Consider a pure property tax whose revenue is rebated equally and whose
model-period rate is $\tau^p(g)$.  With fixed migration, housing clearing and
fiscal balance can be written as
\begin{align}
 F_H(z,g)
&=\log M_H-\log[1-R(z,g)]+\log\bar h(z,g)-\log K(P,g)=0,
 \label{eq:terminal_housing}\\
 F_G(z,g)
 &=T-\tau^p(g)P\bar h(z,g)=0.
 \label{eq:terminal_budget}
\end{align}
The first equation combines stationary population with housing clearing; the
second returns property-tax revenue to households.  At a regular root,
\begin{equation}
 \frac{\dd z^*}{\dd g}=-F_z^{-1}F_g.
 \label{eq:terminal_ift}
\end{equation}
Thus the long-run response can be decomposed into renewal, housing demand,
supply, and fiscal channels even though the derivatives of the Bellman problem
are evaluated numerically.

A fixed-transfer special case makes the decomposition explicit.  Define
\begin{equation}
 \eta_R=-\frac{\partial\log R}{\partial\log P},\quad
 \xi_g=\left.\frac{\partial\log R}{\partial g}\right|_P,\quad
 \eta_h=-\frac{\partial\log\bar h}{\partial\log P},\quad
 d_g=\left.\frac{\partial\log\bar h}{\partial g}\right|_P,
 \label{eq:elasticities}
\end{equation}
and let $\eta_K$ and $s_g$ be the supply elasticity and direct supply shift.
With $\lambda=R/(1-R)$, differentiation of housing clearing yields
\begin{align}
 \frac{\dd\log P^*}{\dd g}
&=\frac{\lambda\xi_g+d_g-s_g}
 {\eta_K+\eta_h+\lambda\eta_R},
 \label{eq:price_elasticity}\\
 \frac{\dd\log S^*}{\dd g}
 &=\lambda\left(\xi_g-\eta_R
 \frac{\dd\log P^*}{\dd g}\right).
 \label{eq:population_elasticity}
\end{align}
The multiplier $\lambda$ captures demographic amplification under fixed
migration.  As renewal approaches one from below, a small change in births per
head has a large effect on population scale; equilibrium prices absorb part of
that response.

In a closed economy, a positive stationary population instead satisfies
$R(P^*,T^*,g)=1$.  Holding the transfer fixed gives
\begin{equation}
 \frac{\dd\log P^*}{\dd g}=\frac{\xi_g}{\eta_R},
 \qquad
 \frac{\dd\log S^*}{\dd g}
 =s_g-d_g+(\eta_K+\eta_h)\frac{\xi_g}{\eta_R}.
 \label{eq:closed_economy_elasticities}
\end{equation}
With an endogenous rebate, the two-equation system analogous to
\eqref{eq:terminal_housing}--\eqref{eq:terminal_budget} should be used.

\section{Transitions and welfare}

For a finite transition, let $z_{0:J}$ stack prices and transfers.  Backward
induction maps this path into values and policies; the household and person
laws map policies and the initial state into future distributions; market and
government clearing produce residuals $F_J(z_{0:J},g)$.  On a branch with
unchanged feasibility and constraint status,
\begin{equation}
 \frac{\dd z_{0:J}}{\dd g}=-F_{J,z}^{-1}F_{J,g}
 \label{eq:transition_ift}
\end{equation}
whenever the path Jacobian is nonsingular.

The forward accounting is exact.  If $\mu_t$ is the household distribution,
$Q_t$ is its transition operator, and $e_t$ contains entrants and externally
specified population flows, then
\begin{align}
 \mu_{t+1}&=Q_t\mu_t+e_t,\\
 \dd\mu_{t+1}&=Q_t\dd\mu_t+(\dd Q_t)\mu_t+\dd e_t.
 \label{eq:distribution_derivative}
\end{align}
The second line separates inherited demographic composition from the response
of current choices.  This propagation is why one permanent change can affect
housing markets for many later dates.

The appropriate reallocation statistic in the quantitative model is
money-metric.  Let $WTP_Y$ be a young household's willingness to pay for access
to a larger housing product, and let $WTA_O$ compensate an old owner for
releasing or downsizing the corresponding unit.  If $\mathcal A_0$ and
$\mathcal A_1$ denote the baseline and expanded access sets, these amounts are
defined by value inversion:
\begin{align}
 V_Y(w_Y-WTP_Y;\mathcal A_1)&=V_Y(w_Y;\mathcal A_0),\\
 V_O(w_O+WTA_O;\text{release})&=V_O(w_O;\text{retain}).
 \label{eq:money_metric_definitions}
\end{align}
If the value comparisons do not already include the real implementation cost
$K$, the surplus is
\begin{equation}
 \mathcal S=WTP_Y-WTA_O-K.
 \label{eq:money_metric_surplus}
\end{equation}
If transaction and moving costs are already included in the two value
comparisons, they must not be subtracted again.  A positive $\mathcal S$ means
that household-specific transfers can produce a Pareto improvement for the
households involved.  It is not a global welfare theorem.

For policy comparisons, a natural first statistic is consumption-equivalent
variation for households alive at the initial date.  A ranking that also
weights households born along different population paths requires an explicit
social criterion.  In addition, the quantitative model does not contain a
realization-based capital-gains tax with stepped-up basis, so it cannot measure
the tax-lock-in terms in the simplified model.

\section{What is established}

The analytical core of the quantitative model is therefore compact.  The
person law gives an exact link from births to renewal and population scale.
The household formulas map housing values into tenure probabilities and birth
hazards.  The equilibrium systems then map those responses into prices,
transfers, and transition paths.  Computing the policy functions, value
differences, and derivatives still requires the numerical model; no global
welfare or transition-uniqueness result follows from these local identities.

\end{document}
